\chapter{Phân tích sâu thuật toán TRACE-CAG}

Đây là chương trọng tâm của báo cáo về mặt kỹ thuật AI. TRACE-CAG là bộ não sinh phản hồi
cho toàn bộ tính năng hội thoại (Lexi Chat, luyện nói, chẩn đoán lỗi) và là điểm khác biệt
chính của LexiLingo so với các ứng dụng học ngôn ngữ khác. Chương~9 đã trình bày công thức
toán học của bộ nhớ đệm SCAR-L1 và hoà trộn truy xuất; chương này trình bày \textbf{toàn bộ
pipeline xung quanh chúng} — từng đỉnh xử lý, thuật toán mở rộng đồ thị, cơ chế thích nghi,
kỹ thuật suy luận đa bước IRCoT, và cách đánh giá chất lượng.

\section{Định nghĩa và động cơ thiết kế}

\textbf{TRACE-CAG} = \en{Trace}-augmented \en{Cache-Augmented Generation}, kết hợp ba ý
tưởng vốn tách rời trong tài liệu học thuật:

\begin{longtable}{L{2.6cm} L{4.2cm} L{6.4cm}}
\toprule
\textbf{Ý tưởng} & \textbf{Nguồn gốc} & \textbf{Vai trò trong TRACE-CAG}\\
\midrule
\endhead
Đồ thị tri thức & GraphRAG & Cấu trúc hoá miền kiến thức tiếng Anh thành khái niệm và quan hệ tiên quyết/liên quan/dễ nhầm, dùng để mở rộng ngữ cảnh có kiểm soát thay vì tìm kiếm mù trên văn bản.\\
Bộ nhớ đệm có xác thực & CAG (Chan et al., 2025) & Tái sử dụng phản hồi đã tính khi trạng thái người học chưa đổi đáng kể, thay vì luôn luôn gọi LLM.\\
Truy vết phụ thuộc & Cung cấp nguồn gốc (\en{provenance}) & Mỗi khoản mục đệm mang theo \en{admissibility certificate} — bằng chứng có thể kiểm chứng để quyết định tái dùng, vá, hay tính lại.\\
\bottomrule
\caption{Ba ý tưởng cấu thành TRACE-CAG}
\end{longtable}

Khác biệt cốt lõi so với RAG kinh điển: RAG nối đất câu trả lời vào \textbf{tài liệu tĩnh}
tìm được qua tìm kiếm véc-tơ; TRACE-CAG nối đất vào \textbf{trạng thái người học đang
sống} (bậc CEFR, lịch sử lỗi, mức nắm vững từng khái niệm) cộng với một đồ thị tri thức
được duyệt có định hướng theo bậc học — không phải tìm kiếm ngữ nghĩa thuần tuý.

\section{Kiến trúc StateGraph}

Pipeline được triển khai bằng \code{StateGraph} của LangGraph: một đồ thị trạng thái nơi
mỗi đỉnh là một hàm bất đồng bộ nhận trạng thái, trả về phần trạng thái cập nhật, và các
cạnh (có thể là cạnh điều kiện) quyết định đỉnh kế tiếp.

\begin{figure}[htbp]
\centering
\resizebox{\textwidth}{!}{%
\begin{tikzpicture}[node distance=7mm and 9mm]
  \node[term] (in) {input\_node};
  \node[dec, right=of in] (dvc) {voice?};
  \node[proc, above right=6mm and 9mm of dvc] (stt) {stt\_node};
  \node[proc, below right=6mm and 9mm of dvc, fill=cSoftO, draw=cGate] (cg) {cache\_gate\_node};
  \node[dec, right=13mm of cg] (dch) {đệm\\trúng?};
  \node[term, above=8mm of dch, fill=cSoftG, draw=cData] (endcache) {END\\(trả từ đệm)};
  \node[proc, below right=6mm and 9mm of dch] (kgd) {kg\_diagnose\_node\\\scriptsize(kg\_expand $\parallel$ diagnose)};
  \node[dec, right=13mm of kgd] (drt) {định tuyến};
  \node[proc, above=9mm of drt] (clar) {ask\_clarify\_node};
  \node[proc, right=13mm of drt] (retr) {retrieve\_node};
  \node[proc, below=9mm of drt] (vn) {vietnamese\_node};
  \node[proc, right=13mm of retr] (gen) {generate\_node};
  \node[term, right=9mm of gen] (end2) {END};

  \draw[flb] (in) -- (dvc);
  \draw[flb] (dvc) -- node[lab]{có} (stt);
  \draw[flb] (dvc) -- node[lab,near start]{không} (cg);
  \draw[flb] (stt) -| (cg);
  \draw[flb] (cg) -- (dch);
  \draw[flb] (dch) -- node[lab]{có} (endcache);
  \draw[flb] (dch) -- node[lab,near start]{không} (kgd);
  \draw[flb] (kgd) -- (drt);
  \draw[flb] (drt) -- node[lab]{$\rho_{diag}{<}0{,}5$} (clar);
  \draw[flb] (drt) -- node[lab,near start]{A1/A2 + lỗi} (vn);
  \draw[flb] (drt) -- node[lab,near start]{khác} (retr);
  \draw[flb] (vn) -| (retr);
  \draw[flb] (retr) -- (gen);
  \draw[flb] (clar) -| (end2);
  \draw[flb] (gen) -- (end2);
\end{tikzpicture}}
\caption{Đồ thị trạng thái TRACE-CAG — chín đỉnh, bốn điểm kết thúc khả dĩ}
\label{fig:tracecag-graph}
\end{figure}

\subsection{Bảng trạng thái pipeline}

\code{TraceCAGState} là một từ điển có kiểu (\code{TypedDict}), chảy qua toàn bộ đồ thị:

\begin{longtable}{L{3.2cm} L{10.2cm}}
\toprule
\textbf{Nhóm} & \textbf{Trường tiêu biểu}\\
\midrule
\endhead
Đầu vào & \code{user\_input}, \code{session\_id}, \code{user\_id}, \code{input\_type}, \code{learner\_profile}, \code{conversation\_history}.\\
Đồ thị tri thức & \code{kg\_seed\_concepts}, \code{kg\_expanded\_nodes}, \code{kg\_paths}, \code{jit\_soft\_graph}.\\
Chẩn đoán & \code{diagnosis\_intent}, \code{diagnosis\_errors}, \code{diagnosis\_root\_causes}, \code{diagnosis\_confidence}.\\
Truy xuất & \code{vector\_hits}, \code{retrieved\_context}, \code{retrieval\_trace}, \code{retrieval\_meta}.\\
Điều khiển thích nghi & \code{adaptive\_profile}, \code{adaptive\_features}, \code{adaptive\_controller}.\\
Điều khiển đệm & \code{cache\_fingerprint}, \code{cache\_decision}, \code{cache\_layer}, \code{cache\_bucket}, \code{reuse\_risk}.\\
Phản hồi & \code{tutor\_response}, \code{native\_hint}, \code{pronunciation\_tip}, \code{strategy}, \code{action\_plan}.\\
Điểm số & \code{fluency\_score}, \code{grammar\_score}, \code{vocabulary\_level}, \code{overall\_score}.\\
\bottomrule
\caption{Nhóm trường trong \code{TraceCAGState}}
\end{longtable}

\section{Đỉnh 1 — \code{input\_node}: nạp ngữ cảnh}

Nạp \code{learner\_profile} từ Redis (khoá \code{profile:\{user\_id\}}, mặc định
\codesp{\{"level": "B1"\}} nếu chưa có) và \code{conversation\_history}
(khoá \code{history:\{session\_id\}}), chuẩn hoá đầu vào. Với đường giọng nói, bản ghi âm
đã được chốt câu bởi lớp STT thời gian thực trước khi vào đỉnh này.

\section{Đỉnh 2 — \code{kg\_expand\_node}: mở rộng đồ thị tri thức}

Đây là thuật toán phức tạp nhất trong toàn pipeline, chia hai giai đoạn.

\subsection{Giai đoạn 1 — Khớp khái niệm hạt giống}

Ba nguồn tín hiệu được hợp nhất để tìm khái niệm hạt giống, theo thứ tự ưu tiên:

\begin{enumerate}[leftmargin=1.4em]
  \item \textbf{Tra cứu từ khoá chính xác} qua chỉ mục đảo ngược — độ phức tạp
  \(O(\text{số từ trong câu})\).
  \item \textbf{Tương đồng ngữ nghĩa TF-IDF} (top 5), bắt được các tham chiếu gián tiếp
  mà tra cứu từ khoá bỏ lỡ.
  \item \textbf{So khớp mẫu lỗi ngữ pháp} bằng biểu thức chính quy — một bảng luật viết tay
  hơn 30 mẫu, mỗi mẫu ánh xạ một cấu trúc câu sai sang một khái niệm ngữ pháp cụ thể.
\end{enumerate}

\begin{longtable}{L{5.4cm} L{6.4cm}}
\toprule
\textbf{Mẫu (rút gọn)} & \textbf{Khái niệm ánh xạ}\\
\midrule
\endhead
\codesp{i goes} & \code{grammar.subject\_verb\_agreement}\\
\codesp{he go} / \codesp{she go} & \code{grammar.third\_person\_s}\\
\codesp{yesterday...(go\textbar want\textbar eat...)} & \code{grammar.past\_time\_markers}\\
\codesp{have went} / \codesp{has came} & \code{grammar.present\_perfect}\\
\codesp{more better} / \codesp{more faster} & \code{grammar.comparatives}\\
\codesp{a apple} / \codesp{a hour} & \code{grammar.articles\_a\_an}\\
\codesp{go to home} / \codesp{married with} & \code{error.preposition\_confusion}\\
\codesp{I very like} & \code{error.word\_order\_svo} (ảnh hưởng SVO tiếng Việt)\\
\codesp{can to} / \codesp{should to} & \code{grammar.modal\_can\_could}\\
\codesp{if...will come} & \code{grammar.conditionals\_first}\\
\codesp{i wish...} / \codesp{if only} & \code{grammar.wish\_regret}\\
\codesp{said that} / \codesp{told...that} & \code{grammar.reported\_speech}\\
\bottomrule
\caption{Trích bảng mẫu lỗi ngữ pháp — nhận diện lỗi đặc trưng của người học Việt Nam}
\label{tbl:grammar-patterns}
\end{longtable}

Ba nguồn được hợp nhất, loại trùng, giữ thứ tự xuất hiện đầu tiên (từ khoá chính xác đứng
trước vì độ tin cậy cao hơn).

\subsection{Giai đoạn 2 — Mở rộng tốt nhất trước có trọng số theo bậc}

Với tập khái niệm hạt giống, thuật toán duyệt đồ thị bằng hàng đợi ưu tiên (giải thuật
tốt nhất trước — \en{best-first search}), ưu tiên khái niệm gần bậc CEFR của người học.

\textbf{Trọng số sư phạm (PedWeight):}
\begin{equation}
w(\ell) =
\begin{cases}
1{,}0 & |\ell - \ell_{\text{người học}}| = 0\\
0{,}7 & |\ell - \ell_{\text{người học}}| = 1\\
0{,}3 & |\ell - \ell_{\text{người học}}| \ge 2
\end{cases}
\end{equation}

với \(\ell\) là thứ tự bậc CEFR (A1=1, \ldots, C2=6). Khái niệm đúng ngay bậc người học
được ưu tiên tuyệt đối; khái niệm quá xa bậc gần như bị bỏ qua nhưng vẫn có mặt trong
hàng đợi để không bỏ sót một quan hệ hiếm nhưng quan trọng.

\begin{lstlisting}[language=Python, caption={Best-first expansion, rút gọn từ \code{kg\_service\_v3.py}}]
frontier = []  # hang doi uu tien: (-w, hop_depth, concept_id, parent, relation)
for s in seed_nodes:
    heappush(frontier, (-1.0, 0, s, None, None))

while frontier and len(expanded_nodes) < max_nodes:      # max_nodes = 10
    neg_w, depth, cid, parent, rel = heappop(frontier)
    if depth > 0:
        expanded_nodes.append(Node(cid, relation=rel, depth=depth))
        if parent: paths.append(Path([parent, cid], [rel]))
    if depth >= max_hops:                                  # max_hops = 2
        continue
    for neighbor_id, edge_rel, neighbor_level in kuzu_neighbors(cid, limit=50):
        if neighbor_id not in visited:
            visited.add(neighbor_id)
            heappush(frontier, (-ped_weight(neighbor_level), depth+1,
                                 neighbor_id, cid, edge_rel))
\end{lstlisting}

\begin{luuy}[Vì sao giới hạn 50 hàng xóm mỗi khái niệm]
Một khái niệm từ vựng phổ biến như "person" hay "run" có thể có 200--500 cạnh ra, trong
khi \code{max\_nodes} chặn toàn bộ lượt duyệt ở khoảng 10. Lấy hết hàng xóm của một
\en{hub node} rồi vứt gần hết sẽ tốn một truy vấn KuzuDB đắt một cách vô ích, nên giới
hạn \codesp{LIMIT 50} cắt chi phí ngay tại nguồn.
\end{luuy}

Kết quả được đệm ở Redis theo khoá
\code{kg:subgraph:sha1(sorted(seeds):level:max\_hops:max\_nodes)} với thời gian sống
30 phút, để các yêu cầu có cùng hạt giống và cùng bậc không phải duyệt lại đồ thị.

\section{Đỉnh 3 — \code{diagnose\_node}: chẩn đoán bằng LLM}

Chạy \textbf{song song} với \code{kg\_expand\_node} qua \code{asyncio.gather} — hai việc
độc lập dữ liệu đầu vào nên chạy đồng thời cắt gần một nửa độ trễ nhánh chậm. Lời nhắc gửi
LLM có cấu trúc:

\begin{lstlisting}
"{cau cua nguoi hoc} [learner: level={bac}, native={tieng_me_de}]"
"Identify: intent, grammar errors, fluency issues, confidence"
\end{lstlisting}

Kỳ vọng đầu ra là JSON có cấu trúc:

\begin{lstlisting}
{
  "intent": "correct",             // correct | explain | practice | ask
  "errors": [
    {"span": "goes", "type": "verb_conjugation",
     "correction": "went", "explanation": "Past simple irregular verb"}
  ],
  "root_cause_concepts": ["past_simple", "irregular_verbs"],
  "confidence": 0.92,
  "fluency_score": 0.75, "grammar_score": 0.65
}
\end{lstlisting}

Nếu LLM lỗi hoặc trả về không đúng định dạng, \code{\_rule\_based\_diagnosis} là đường dự
phòng dựa trên chính bảng mẫu regex ở mục~\ref{tbl:grammar-patterns} — nghĩa là hệ thống
không bao giờ hoàn toàn im lặng ngay cả khi mọi nhà cung cấp LLM đều gặp sự cố.

\section{Định tuyến sau chẩn đoán}

\begin{lstlisting}[language=Python]
def route_after_diagnosis(state):
    confidence = state.get("diagnosis_confidence", 1.0)
    level      = state["learner_profile"].get("level", "B1")
    errors     = state.get("diagnosis_errors", [])
    if confidence < 0.5:
        return "ask_clarify_node"        # khong doan mo
    if level in ("A1", "A2") and errors:
        return "vietnamese_node"         # giai thich tieng me de truoc
    return "retrieve_node"
\end{lstlisting}

\code{ask\_clarify\_node} đưa ra ba lựa chọn (Sửa lỗi / Giải thích / Luyện tập) và kết
thúc luôn — không lãng phí một lượt gọi LLM sinh phản hồi khi hệ thống chưa đủ tin cậy về
điều người học thực sự muốn.

\section{Đỉnh 4 — \code{retrieve\_node}: truy xuất lai có ngân sách}

Đây là điểm khác biệt lớn thứ hai so với RAG thường: truy xuất bị \textbf{giới hạn ngân
sách thời gian theo từng giai đoạn}, không chạy tới khi xong mới dừng.

\begin{longtable}{L{4.2cm} L{2.6cm} L{6.2cm}}
\toprule
\textbf{Giai đoạn} & \textbf{Ngân sách mặc định} & \textbf{Nội dung}\\
\midrule
\endhead
Đồ thị (KG) & 120 ms & Lắp ráp ngữ cảnh từ \code{kg\_expanded\_nodes} và \code{jit\_soft\_graph}.\\
Véc-tơ & 80 ms & \code{RetrievalServiceV3} (xếp hạng theo độ trung tâm và cộng đồng đồ thị), dự phòng bằng cổng MiniLM nếu lỗi.\\
Hoà trộn & 40 ms & Tính điểm hợp nhất (công thức Chương~9) và xếp hạng lại bằng bộ xếp hạng học trực tuyến.\\
\bottomrule
\caption{Ngân sách ba giai đoạn của truy xuất lai}
\end{longtable}

Giai đoạn véc-tơ được khởi chạy bằng \code{asyncio.create\_task} \textbf{song song} với
giai đoạn đồ thị, vì đầu vào của nó (\code{kg\_concepts}, \code{user\_input}) không phụ
thuộc kết quả giai đoạn đồ thị. Nếu hết ngân sách, cờ \code{budget\_exhausted} được bật và
pipeline đi tiếp với những gì đã có — \textbf{trễ có giới hạn trên} luôn được ưu tiên hơn
kết quả đầy đủ nhưng không đoán trước được thời gian.

Công thức hoà trộn ba tín hiệu (đồ thị, véc-tơ, độ mới) và mười hai đặc trưng của bộ xếp
hạng học trực tuyến đã trình bày đầy đủ ở mục~6.5.

\section{Đỉnh 5 — \code{generate\_node}: sinh phản hồi có chiến lược}

\subsection{Chọn chiến lược sư phạm}

\begin{longtable}{L{3.4cm} L{3.6cm} L{5.6cm}}
\toprule
\textbf{Điều kiện} & \textbf{Chiến lược} & \textbf{Hành vi}\\
\midrule
\endhead
Ý định \code{explain} và độ trôi chảy $>0{,}7$ & \en{socratic} & Dẫn dắt bằng câu hỏi ngắn, không lộ đáp án ngay; chỉ đưa lỗi dưới dạng gợi ý.\\
Không có lỗi & \en{praise} & Khen và khuyến khích tiếp tục.\\
1--2 lỗi & \en{feedback} & Nêu lỗi kèm giải thích ngắn gọn, tự nhiên.\\
Từ 3 lỗi trở lên & \en{scaffold} & Chia nhỏ, hỗ trợ từng bước — tránh làm người học quá tải thông tin.\\
\bottomrule
\caption{Bốn chiến lược sư phạm và điều kiện kích hoạt}
\end{longtable}

\subsection{Độ dốc độ khó theo tiến trình phiên}

\begin{equation}
\text{độ khó} =
\begin{cases}
\text{gentle} & \text{lượt} < 3 \text{ hoặc điểm tổng} < 0{,}50\\
\text{standard} & \text{lượt} < 8 \text{ hoặc điểm tổng} < 0{,}70\\
\text{challenging} & \text{ngược lại}
\end{cases}
\end{equation}

Độ khó này không quyết định \textit{nội dung} mà quyết định \textit{cách trình bày} — đưa
vào lời nhắc hệ thống để LLM điều chỉnh độ phức tạp câu chữ, có nương nhẹ ban đầu và siết
dần theo tiến trình phiên hội thoại.

\subsection{Cấu trúc lời nhắc hệ thống}

\begin{lstlisting}
You are LexiLingo, an English tutor.
Learner: level={bac}, native={tieng_me_de}, difficulty={do_kho}

--- Knowledge Graph Context ---
{retrieved_context}

Strategy: {strategy}.
--- Errors Found ---           (neu co, khong dung cho socratic)
- 'goes' -> 'went' (Past simple irregular verb)
\end{lstlisting}

Sau khi LLM sinh xong: trích \code{tutor\_response}, cập nhật lịch sử hội thoại, cập nhật
điểm mức nắm vững trên đồ thị, và nếu \codesp{cache\_decision == "full"} thì ghi khoản mục
đệm mới với chứng chỉ khả dụng (Chương~9).

\section{Điều khiển thích nghi — lựa chọn cấu hình bằng bandit ràng buộc}

Khi \code{retrieval\_policy} là \code{"adaptive"}, hệ thống không dùng một bộ ngưỡng cố
định cho mọi yêu cầu mà chọn giữa ba \textbf{cấu hình} (\en{profile}), mỗi cấu hình đánh
đổi tốc độ lấy chất lượng theo mức khác nhau:

\begin{longtable}{L{2.0cm} L{2.0cm} L{2.0cm} L{2.0cm} L{2.4cm}}
\toprule
\textbf{Cấu hình} & \(\Delta\tau_{\text{reuse}}\) & \(\Delta\tau_{\text{patch}}\) & \textbf{Hệ số chất lượng} & \textbf{Chi phí độ trễ}\\
\midrule
\endhead
fast & $-0{,}06$ & $-0{,}04$ & 0,86 & 0,40\\
balanced & 0,00 & 0,00 & 1,00 & 0,62\\
quality & $+0{,}04$ & $+0{,}08$ & 1,06 & 0,88\\
\bottomrule
\caption{Ba cấu hình thích nghi — cấu hình \en{fast} nới ngưỡng tái dùng đệm để ưu tiên tốc độ}
\end{longtable}

\subsection{Hàm mục tiêu và lựa chọn}

Với mỗi yêu cầu, hệ thống tính hàm mục tiêu cho từng cấu hình dựa trên đặc trưng ngữ cảnh
(áp lực nhà cung cấp LLM, độ dài câu, số lượt hội thoại) cộng với độ lệch học được
(\en{bias}) từ phản hồi các lần trước, rồi chọn cấu hình có mục tiêu cao nhất — trừ
5\% xác suất khám phá ngẫu nhiên (\(\epsilon = 0{,}03\)) để tránh kẹt cứng vào một lựa
chọn cục bộ tối ưu.

\subsection{Cập nhật trực tuyến có ràng buộc kép}

Sau khi có kết quả thật (đúng cho lô \en{benchmark}), phần thưởng được tính theo thứ tự
ưu tiên sản phẩm yêu cầu: \textbf{chất lượng trước, truy xuất sau, an toàn trôi dạt, rồi
mới đến độ trễ}:

\begin{equation}
r_{\text{thô}} = 0{,}46\cdot F_1 + 0{,}16\cdot EM + 0{,}24\cdot(0{,}5R@1{+}0{,}3R@3{+}0{,}2R@5) + 0{,}14\cdot\mathbb{1}[\text{tái dùng an toàn}]
\end{equation}

Hai \textbf{biến đối ngẫu} (\en{dual variables}) phạt việc vượt mục tiêu độ trễ
(380 ms) và tỉ lệ tái dùng sai (12\%):

\begin{equation}
r = r_{\text{thô}} - \lambda_{\text{trễ}}\cdot\max(0, \Delta_{\text{trễ}}) - \lambda_{\text{sai}}\cdot\max(0, \Delta_{\text{sai}})
\end{equation}

Độ lệch của cấu hình được chọn dịch chuyển theo phần thưởng \textbf{đã trừ trung bình
chạy} (giảm phương sai), và hai biến đối ngẫu tự tăng khi mục tiêu ràng buộc bị vi phạm —
đây là một bộ điều khiển kiểu \en{primal-dual} cỡ nhỏ, chạy hoàn toàn trong bộ nhớ tiến
trình mà không cần huấn luyện ngoại tuyến.

\begin{luuy}[Chỉ dùng trong bối cảnh đánh giá có kiểm soát]
Cơ chế thích nghi hiện chỉ bật khi \code{retrieval\_policy="adaptive"} hoặc trong chế độ
\en{benchmark}. Nó chưa phải đường mặc định cho lưu lượng người dùng thật.
\end{luuy}

\section{IRCoT — suy luận rồi truy xuất lại cho câu hỏi đa bước}

Với những câu hỏi cần bắc cầu qua hai thực thể (ví dụ: \textit{"đạo diễn của bộ phim mà
diễn viên X đóng chính là ai?"}), truy xuất một lượt thường bỏ lỡ thực thể trung gian. Cơ
chế \textbf{IRCoT} (\en{Interleaving Retrieval with Chain-of-Thought}, arXiv:2212.10509)
giải quyết bằng một vòng suy luận $\rightarrow$ truy xuất lại có định hướng:

\begin{enumerate}[leftmargin=1.4em]
  \item \textbf{Suy luận}: hỏi LLM "thực thể bắc cầu nào cần tra tiếp theo?", chỉ nhận về
  một dòng tên thực thể duy nhất.
  \item \textbf{Truy xuất lại có định hướng}: dùng thực thể đó tìm tối đa hai đoạn văn mới
  từ toàn bộ kho ứng viên (không chỉ những gì đã có trong ngữ cảnh).
  \item \textbf{Xác thực hợp đồng}: kiểm tra các đoạn mới tìm được có thực sự liên quan
  câu hỏi và thực thể bắc cầu hay không trước khi chấp nhận; nếu không đạt, giữ nguyên
  ngữ cảnh gốc.
  \item \textbf{Hợp nhất}: đoạn bắc cầu được đặt \textbf{trước} ngữ cảnh gốc, dùng ngân
  sách ký tự riêng lớn hơn ngân sách nền tảng để tránh đoạn chứa câu trả lời gốc bị đẩy
  ra ngoài khi cắt bớt.
\end{enumerate}

\begin{luuy}[Khác biệt so với truy xuất một lượt]
Đây là cơ chế lặp \textbf{suy luận$\rightarrow$truy xuất lại$\rightarrow$trả lời} đúng
nguyên bản của IRCoT, khác với các thử nghiệm một lượt (mở rộng đồ thị tĩnh, hay chỉ tăng
số lượng đoạn lấy về): bước suy luận \textbf{chủ động dẫn dắt} một lượt truy xuất tập
trung vào đúng đoạn văn bước hai mà xếp hạng từ vựng thông thường bỏ sót do giới hạn ngân
sách.
\end{luuy}

Trên bộ đánh giá đa bước nội bộ, việc thêm IRCoT sau khi sửa hai lỗi tự gây ra (ngữ cảnh
bị đẩy ra do cắt bớt, và gợi ý định hướng quá chặt tay) đã cải thiện \textbf{Exact Match}
thêm 3,1 điểm phần trăm, xác nhận rằng nút thắt nằm ở khâu \textbf{tổng hợp suy luận đa
bước của bộ đọc}, không phải ở kích thước mô hình hay chất lượng truy xuất đơn bước.

\section{Định dạng phản hồi cuối cùng}

Điểm số luôn được \textbf{tính lại} tại \code{\_format\_response} bất kể pipeline đi
nhánh nào (đệm hay tính mới), để hai đường này không bao giờ trả về ngữ nghĩa điểm số
khác nhau.

\begin{lstlisting}
{
  "tutor_response": "...",
  "corrections":    [ {error, correction, type, explanation} ],
  "linked_concepts": ["past_simple", "irregular_verbs"],   // tu khoa cau hoi
  "weak_concepts":   ["past_simple"],                       // loi that su, de goi y
  "scores": {
    "fluency": 0.75, "grammar": 0.65, "overall": 0.69,
    "retrieval": {"precision_k": 0.8, "recall_k": 0.6, "ndcg_k": 0.7512, "mrr": 0.6667}
  },
  "action":   {"strategy": "scaffold", "next": "continue"},
  "metadata": {"latency_ms": 1250, "ttft_ms": 340, "cache_hit": false,
               "cache_decision": "full", "reuse_risk": 1.0,
               "models_used": ["gemini-2.0-flash"], "path": "slow"}
}
\end{lstlisting}

\section{Đánh giá chất lượng}

\subsection{Điểm tổng hợp}

\begin{equation}
\text{overall} = 0{,}40\cdot\text{grammar} + 0{,}30\cdot\text{fluency} + 0{,}30\cdot v(\text{vocab\_level})
\end{equation}

thay cho công thức \(0{,}6\cdot\text{grammar}+0{,}4\cdot\text{fluency}\) ngây thơ trước đây
— vốn hoàn toàn bỏ qua tín hiệu từ vựng.

\subsection{Chỉ số chất lượng văn bản và truy xuất}

WER, TTR, mật độ lỗi, Precision@k, Recall@k, NDCG@k, MRR đã trình bày công thức đầy đủ ở
mục~6.6. \code{EvaluationAgent} tính lại các chỉ số này cho \textbf{mọi} phản hồi, kể cả
lấy từ đệm, để phục vụ theo dõi chất lượng theo thời gian một cách nhất quán.

\subsection{Hiệu năng theo kịch bản}

\begin{longtable}{L{4.4cm} L{2.0cm} L{2.4cm} L{4.0cm}}
\toprule
\textbf{Kịch bản} & \textbf{Nhánh} & \textbf{Độ trễ} & \textbf{Mô hình}\\
\midrule
\endhead
Trúng đệm L0 (chính xác) & nhanh & $<10$ ms & không gọi\\
Trúng đệm L1 (theo \en{bucket}) & nhanh & $<50$ ms & không gọi\\
Câu chào đơn giản & chậm & $\approx$3 s & mô hình nhỏ cục bộ\\
Câu hỏi ngữ pháp & chậm & $\approx$2 s & mô hình đám mây\\
Phân tích phức tạp & chậm & $\approx$20 s & mô hình suy luận cục bộ\\
Giọng nói đầy đủ & chậm & 5--8 s & STT + LLM + phân tích âm vị\\
\bottomrule
\caption{Đặc tính hiệu năng theo kịch bản sử dụng}
\end{longtable}

\subsection{Kết quả benchmark HotpotQA (n=64, \code{qwen/qwen3.6-27b})}

Đường cơ sở HippoRAG trong các lần đo trước là bản \textbf{tự cài lại}
(\code{hipporag\_proxy}), không đủ tư cách chứng minh vượt trội so với công trình
đã công bố. Từ 2026-08-26, gói HippoRAG~2 chính thức (\code{pip install hipporag})
được chạy như đường cơ sở ngoài, cùng bộ dữ liệu, cùng mô hình, chấm bằng đúng
hàm chuẩn hoá SQuAD của repo.

\begin{longtable}{L{3.8cm} R{1.8cm} R{1.8cm} R{2.0cm} R{2.2cm}}
\toprule
\textbf{Hệ thống} & \textbf{EM} & \textbf{F1} & \textbf{R@5} & \textbf{Độ trễ}\\
\midrule
\endhead
CAG thuần & 62,5\% & 74,5\% & 74,2\% & 2,1 s\\
\textbf{TRACE-CAG} & \textbf{62,5\%} & \textbf{74,7\%} & 78,9\% & 2,2 s\\
HippoRAG~2 (bản chính thức) & 53,1--56,2\% & 69,7--73,0\% & \textbf{85,9\%} & 66--81 s\\
\bottomrule
\caption{So sánh trên HotpotQA, n=64, cùng mô hình đọc (cập nhật 2026-08-29).
Cột thứ tư là R@5. Độ trễ là trung bình lượt lạnh; HippoRAG lập chỉ mục lại kho
10 đoạn cho mỗi câu. Số EM/F1 của HippoRAG là khoảng qua hai lần chạy cùng cấu
hình --- 4/64 câu khác nhau do nhiễu mô hình, không phải do cấu hình.}
\end{longtable}

Bảng này thay cho bảng cũ và \textbf{thu hẹp kết luận}. HippoRAG~2 chạy đủ 64/64 với
ngân sách sinh đáp án 1024 token thực sự có hiệu lực --- trước đó phải đặt
\code{max\_completion\_tokens} chứ không phải \code{max\_tokens}, nên tham số bị
thư viện ghi đè âm thầm và hai lần chạy đầu thực chất cùng một cấu hình. Nhờ vậy
số của HippoRAG tăng lên 56,2\%/73,0\%.

TRACE-CAG vẫn vượt gói HippoRAG chính thức ở cả hai chỉ số chính (hơn 6,3--9,4
điểm EM) với độ trễ thấp hơn $\approx$30--37 lần. Nhưng sau khi sửa lỗi gán nhãn của bộ xếp hạng và lỗi IRCoT
(2026-08-28), \textbf{TRACE-CAG không còn vượt CAG thuần về EM}: hai bên bằng nhau,
và chênh lệch F1 0,4 điểm nằm trong nhiễu giữa các lần chạy. Các bản sửa đó mang
tính chung nên nâng đường cơ sở nhiều hơn nâng nhánh đồ thị. Lợi thế truy xuất thì
vẫn còn (\code{all\_support@5} 59,4\% so với 50,0\%), nhưng \textbf{không còn chuyển
hoá thành chất lượng câu trả lời} --- nút thắt đã dời sang bộ đọc.

\textbf{Chỉ số R@5 của HippoRAG đã được sửa và đo lại.} Bộ đo cũ tìm tiêu đề vàng
trong \emph{toàn văn gộp} của các đoạn lấy về; trên câu bắc cầu HotpotQA thì thân
đoạn này gọi tên tiêu đề đoạn kia theo đúng định nghĩa, nên nó tính công cho cả
những đoạn chưa hề được truy hồi. Sau khi sửa về so khớp danh tính tài liệu và
chạy lại đủ n=64 (0 lỗi): R@5 giảm từ 93,8\% xuống \textbf{85,9\%}, tức bộ đo cũ
thổi phồng \textbf{7,8 điểm} và lệch ở 13/64 câu. Đáng chú ý nó nhiễu \emph{hai
chiều} chứ không chỉ rộng rãi một chiều: ở 2 câu, đoạn vàng có được truy hồi
nhưng tiêu đề của nó không xuất hiện trong phần thân nào, nên bộ đo cũ lại chấm
trượt một lần trúng.

Một ước lượng trước đó trong tài liệu này ghi ``46,9\% tiêu đề vàng có thể trích
từ đoạn khác''. Đó là \emph{cơ hội} thổi phồng chứ không phải mức thổi phồng, và
nó phóng đại khoảng 6 lần. Bài học: đo chỉ số bằng một lần chạy lại, đừng suy ra
từ đặc điểm của kho ngữ liệu.

\begin{luuy}[Nút thắt thực sự là khả năng với tới bằng chứng, không phải suy luận]
Ba chỉ số truy xuất được bổ sung vì không chỉ số cũ nào phân biệt được ``tìm thấy
một bước'' với ``tìm thấy câu trả lời'': \code{hit@k} (có ít nhất một đoạn hỗ trợ),
\code{all\_support@k} (có \textbf{đủ} mọi đoạn hỗ trợ) và \code{answer\_in\_context@k}
(chuỗi đáp án nằm trong cửa sổ mà mô hình đọc thấy --- trần cứng của bộ đọc).

Đo được \code{hit@5}$\approx$100\% trong khi \code{all\_support@5}=57,8\%: một bước
hầu như luôn tới, cả hai thì thường không. Khoảng cách này vô hình với
\code{recall@k} vì recall lấy trung bình trên các bước, nên chấm 0,5 cho một câu
thực chất không thể trả lời được. Tỉ lệ chuyển hoá của bộ đọc
(EM chia \code{answer\_in\_context}) khi đó là \textbf{90\%} --- bộ đọc đã khai thác
gần hết những gì nó nhìn thấy.

Nâng ngân sách bằng chứng từ 5 lên 9--10 đoạn đưa \code{answer\_in\_context} từ
62,5\% lên 82,8\% và EM từ 56,2\% lên \textbf{62,5\%} (chênh lệch được giải thích
trọn vẹn ở mức từng câu: 6 câu thắng, 2 câu thua), đổi lại 29\% token đầu vào.

Kết quả này \textbf{bác bỏ} kết luận trước đó rằng ``suy luận đa bước của bộ đọc mới
là ràng buộc''. Lần thử trước nâng \emph{trọng số} đồ thị --- tức sắp xếp lại cùng
một ngân sách cố định --- và EM giảm; nâng \emph{ngân sách} lại làm EM tăng mạnh.
Sắp xếp lại và mở rộng khả năng với tới là hai đòn bẩy khác nhau, không được suy
diễn kết quả của cái này sang cái kia.
\end{luuy}

\begin{luuy}[Cập nhật 2026-08-29: nút thắt đã dời sang bộ đọc, và ba kỹ thuật bị bác bỏ]
Con số ``chuyển hoá 90\%'' ở trên thuộc về điểm vận hành cũ (ngân sách 5 đoạn).
Ở ngân sách 9--10 đoạn hiện tại, \code{answer\_in\_context} là 82,8\% còn EM là
62,5\%, tức chuyển hoá \textbf{75,5\%}: gần một phần tư số câu có sẵn đáp án trong
ngữ cảnh mà bộ đọc vẫn trả sai. Dư địa còn lại là 20,3 điểm và nó nằm ở bộ đọc.

Phân loại 24 câu sai (\code{benchmark/analyze\_failures.py}): 10 câu sai độ chi tiết
của span, 7 câu chọn nhầm span, 4 câu đáp án không có trong ngữ cảnh, 2 câu bị hậu
xử lý bác bỏ, 1 câu đúng/sai. Quan trọng nhất: \textbf{18 trong 24 câu sai có đáp án
là span nguyên văn trong ngữ cảnh} --- mô hình chọn nhầm span chứ không cắt sai biên.

Ba kỹ thuật đã đo và \textbf{bác bỏ}, ghi lại để không ai thử lại:
\begin{itemize}
  \item \emph{Ép về span}: lỗi hai chiều (đáp án vàng khi dài hơn, khi ngắn hơn dự
        đoán) và nhiều đáp án là số hoặc chữ thường, nên luật dựa vào chữ hoa không
        phủ được. Khớp với lần thử trước: 0 câu thắng, 3--4 câu thua.
  \item \emph{Cò ``không phải span nguyên văn''}: chỉ 46\% chính xác --- đụng vào 7
        câu đang đúng mà chỉ với tới 6 trong 24 câu sai.
  \item \emph{Tự nhất quán} (5 mẫu, nhiệt độ 0,7): 0 câu thắng, 1 câu thua. Tỉ lệ đổi
        đáp án chỉ 2/85: với 4--5 phiếu, đáp án tham lam thắng 83/85 lần. Chi phí
        175 giây mỗi câu.
\end{itemize}

Kết luận: \textbf{lỗi của bộ đọc là có hệ thống, không ngẫu nhiên}. Mô hình chọn cùng
một span sai một cách nhất quán, nên mọi kỹ thuật dựa trên lấy mẫu và bỏ phiếu đều
không sửa được. Phần dư địa còn lại đòi hỏi bộ đọc mạnh hơn hoặc giám sát ở mức span,
không phải một lớp bọc quanh cùng mô hình.
\end{luuy}

Chín lỗi kỹ thuật đã được sửa trong quá trình đánh giá (đặt tên sai chế độ, cắt bớt tập dữ
liệu, ngân sách bằng chứng, cắt ngữ cảnh, khoá Groq lỗi, không thử lại khi gặp lỗi 503,
mốc thời gian tạo không đơn điệu qua Redis, chuẩn hoá SQuAD không đúng chuẩn, phân loại sai
mô hình mạnh/yếu), cùng ba lỗi thuật toán tìm được ở vòng 2026-08-26: bộ trích xuất
\en{anchor} dán từ để hỏi vào thực thể đầu tiên (khiến một bước không có \en{anchor} nào
dùng được), chọn bằng chứng thưởng độ phủ tuyệt đối thay vì độ phủ \textbf{tăng thêm}, và
cổng chọn lọc IRCoT có ba nhánh chết vì đọc một trường mà thời gian chạy không bao giờ
truyền vào. Bốn lần thử tối ưu khác (trọng số đồ thị, đơn giản hoá E2G, đổi sang mô hình
70B) đều \textbf{làm giảm} chất lượng và bị hoàn tác.

\begin{luuy}[Giới hạn cần nêu khi công bố]
Kết quả dựa trên một lát cắt n=64 duy nhất. Phần lớn mức tăng đến từ ngân sách bằng
chứng --- một \textbf{tham số giao thức đánh giá}, không phải đóng góp kiến trúc.

Con số \textbf{+3,1 EM / +3,4 F1} từng được ghi ở đây như ``đóng góp kiến trúc riêng''
\textbf{không còn đúng} và đã bị gỡ: sau các bản sửa 2026-08-28, khoảng cách so với CAG
thuần là \textbf{0,0 EM / +0,3 F1}. Tương tự, tỉ lệ trúng đệm 48,4\% không phải một đóng
góp kiến trúc --- xem mục~\ref{sec:tracecag-cache-decomp}: toàn bộ là đệm trùng khớp
chính xác ở lượt ấm, còn tầng L1 theo bucket đồ thị chưa trúng lần nào.

Ngoài ra \code{answer\_in\_context} là 82,8\%, nghĩa là $\approx$17\% câu hỏi nằm ngoài
tầm với ngay cả khi lấy toàn bộ 10 đoạn --- không thay đổi nào ở bộ đọc cứu được nhóm này.
\end{luuy}

\section{Phân rã chỉ số: điều bảng tổng không cho thấy}
\label{sec:tracecag-cache-decomp}

Bảng so sánh ở trên trả lời ``hệ nào cao điểm hơn''. Ba phép phân rã dưới đây trả lời câu
hỏi khó hơn --- \emph{vì sao} --- và cả ba đều thu hẹp tuyên bố. Số liệu lấy từ cùng một
lần chạy \code{hotpotqa\_n64\_cleanup} cho hai chế độ nội bộ.

\subsection{TRACE-CAG thắng mọi chỉ số truy xuất và không thắng chỉ số đáp án}

\begin{longtable}{L{4.6cm} R{2.6cm} R{2.6cm} R{2.0cm}}
\toprule
\textbf{Chỉ số} & \textbf{CAG thuần} & \textbf{TRACE-CAG} & \textbf{$\Delta$}\\
\midrule
\endhead
EM & 62,5\% & 62,5\% & \textbf{$\pm$0,0}\\
F1 & 74,5\% & 74,7\% & +0,3\\
\midrule
recall@5 & 74,2\% & 78,9\% & +4,7\\
precision@5 & 30,6\% & 32,5\% & +1,9\\
\code{all\_support@5} & 50,0\% & 59,4\% & \textbf{+9,4}\\
\code{answer\_in\_context@5} & 57,8\% & 62,5\% & +4,7\\
MRR@5 & 84,4\% & 87,0\% & +2,6\\
nDCG@5 & 71,3\% & 74,7\% & +3,4\\
MAP@5 & 63,3\% & 66,4\% & +3,1\\
\bottomrule
\caption{Hai chế độ chạy trong cùng một tiến trình benchmark, cùng seed, cùng nhà cung
cấp; khác nhau đúng một thứ là nhánh đồ thị. Đơn vị $\Delta$ là điểm phần trăm.}
\end{longtable}

Nhánh đồ thị đưa thêm 9,4 điểm số câu vào trạng thái ``đủ cả hai đoạn hỗ trợ'' --- điều
kiện cần để trả lời một câu hai bước --- mà EM không nhúc nhích. Không thể đổ cho khâu cắt
ngữ cảnh, vì \code{answer\_in\_context@5} cũng tăng 4,7 điểm: chuỗi đáp án \emph{thật sự}
lọt vào cửa sổ mô hình nhìn thấy thường xuyên hơn.

Đối sánh từng câu ở lượt lạnh cho thấy lý do trực tiếp: \textbf{hai chế độ trả lời giống hệt
nhau ở 59 trong 64 câu} (92\%). Chỉ 5 câu khác chuỗi, và 4 câu trong đó đổi điểm --- chia
đều hai chiều, hai câu mỗi bên. EM bằng nhau ở đây không phải hai hệ khác nhau tình cờ trùng
điểm: ngữ cảnh tốt hơn đi vào bộ đọc, cùng một câu trả lời đi ra.

\subsection{Tỉ lệ trúng đệm 48,4\% đo cái gì}

Giao thức đặt \code{cache\_repeats: 2}: mỗi câu chạy hai lượt, một lạnh một ấm (64 câu
$\to$ 128 quan sát mỗi chế độ). Phân rã theo tầng đệm trên toàn bộ 128 quan sát:

\begin{longtable}{L{5.6cm} R{3.0cm} R{3.0cm}}
\toprule
\textbf{Tầng} & \textbf{CAG thuần} & \textbf{TRACE-CAG}\\
\midrule
\endhead
L0 --- trùng khớp chính xác & 61 & 62\\
\textbf{L1 --- theo bucket đồ thị} & \textbf{0} & \textbf{0}\\
không trúng & 67 & 66\\
\bottomrule
\caption{Trúng đệm ở lượt lạnh bằng 0 cho cả hai chế độ; toàn bộ số trúng nằm ở lượt ấm.}
\end{longtable}

Nên 48,4\% \textbf{không phải tỉ lệ tái dùng ngữ nghĩa}. Nó là trần 50\% do giao thức áp
đặt, và TRACE-CAG chạm 96,9\% của trần đó. Điều đo được là ``hỏi lại đúng câu vừa hỏi thì
trả lại được'' --- một tính chất của L0, mà bất kỳ hệ nào cũng cài được.

Đây là điểm phải nói thẳng trong mọi bản công bố: \textbf{L1 mới là cơ chế phân biệt
TRACE-CAG với một CAG có đệm thường}, và \code{l1\_rate} bằng 0 ở mọi chế độ của mọi lần
chạy kể từ 2026-05-30.

\subsection{Vì sao L1 bằng 0: đường ghi bị đóng, không phải thiếu dữ liệu}

Lời giải thích lưu truyền từ 2026-05-31 là ``benchmark dùng mẫu độc nhất nên không có trùng
lặp ngữ nghĩa; L1 sẽ hiệu quả trong production''. Lời giải thích đó \textbf{sai}, và bằng
chứng phản bác nằm sẵn trong repo: bộ dữ liệu \code{query\_clusters} được tạo đúng để kiểm
L1 --- gồm các cụm câu hỏi diễn đạt lại của cùng một câu --- đã chạy ngày 2026-07-10 với
validation PASS, và vẫn cho \code{l1\_rate} bằng 0.

Truy vết trong dữ liệu chạy: cả 64/64 quan sát lạnh kết thúc với
\code{cache\_gate\_meta.reasons} bằng \code{('no\_compatible\_candidate',)}. Trong
\code{cache\_gate\_node}, đó là nhánh mặc định khi biến \code{best\_rejection} vẫn là
\code{None} --- nghĩa là \textbf{vòng lặp duyệt ứng viên chưa chạy dù chỉ một lần}. Không có
ứng viên nào để mà từ chối; vấn đề nằm \emph{trước} cổng PCC chứ không phải tại cổng PCC.

Ngược lên khâu ghi: \code{\_register\_l1\_bucket\_aliases} chỉ được gọi khi
\code{\_is\_pcc\_stable(state)} đúng, và hàm đó chấp nhận theo một trong hai lối --- cả hai
đều đóng trong benchmark:

\begin{longtable}{L{3.0cm} L{5.4cm} L{5.6cm}}
\toprule
\textbf{Lối} & \textbf{Điều kiện} & \textbf{Benchmark cấp gì}\\
\midrule
\endhead
oracle & \code{\_tracecag\_state["concepts"]} khác rỗng & \code{\_public\_qa\_state()} chỉ
trả \code{evidence\_hash}, \code{source\_version}, \code{freshness\_class} --- \textbf{không
có \code{concepts}}\\
chẩn đoán & \code{diagnosis\_confidence} $\geq$ 0,70 \textbf{và}
\code{diagnosis\_root\_causes} khác rỗng & nhánh bypass đặt \code{confidence} bằng 1,0 nhưng
\code{root\_causes} bằng rỗng\\
\bottomrule
\caption{Hai lối mở cho việc ghi L1, và lý do benchmark không đi được lối nào. Kiểm chứng
bằng thực thi trực tiếp: \code{\_is\_pcc\_stable} trả \code{False} trên trạng thái đúng như
harness tạo ra, và trả \code{True} khi thêm một trong hai khoá còn thiếu.}
\end{longtable}

Nên \textbf{L1 không thể trúng trong benchmark theo đúng cấu tạo}: không bucket nào từng
được ghi, nên không ứng viên nào để đọc. Đó là lý do \code{query\_clusters} cũng cho 0\%, và
là lý do mọi kế hoạch kiểu ``chạy thêm một tập dữ liệu khác để đo L1'' đều sẽ cho 0\%.

Cần nói rõ sắc thái: tài liệu nội hàm của \code{\_is\_pcc\_stable} \emph{mô tả đúng} ý định
--- nhánh bypass bỏ qua chẩn đoán, và ``khái niệm do oracle cung cấp sẽ neo bucket thay
thế''. Lối thoát oracle đã được thiết kế; chỉ là harness không bao giờ cấp khoá
\code{concepts}. Đây là lệch pha giữa hai module, không phải thiếu sót thiết kế.

\subsection{Sửa và đo: L1 chạy lần đầu tiên}

Ba thay đổi, mỗi thay đổi đo riêng trên \code{query\_clusters} (32 mẫu, một lượt mỗi câu, để
câu diễn đạt lại buộc phải trúng L1 chứ không phải L0).

\begin{longtable}{L{6.2cm} R{2.2cm} R{2.2cm} R{2.2cm}}
\toprule
\textbf{Giai đoạn} & \textbf{EM} & \textbf{F1} & \textbf{L1}\\
\midrule
\endhead
1. mở cổng ghi L1 & 50,0\% & 68,4\% & \textbf{0,0\%}\\
2. + khai báo \code{patchable\_slots} & 25,0\% & 52,3\% & \textbf{75,0\%}\\
3. + lọc token nội bộ khỏi bản vá & \textbf{50,0\%} & 67,9\% & \textbf{75,0\%}\\
\bottomrule
\caption{Cả ba lần đều 0 lỗi, validation PASS.}
\end{longtable}

\textbf{Giai đoạn 1} cho \code{\_is\_pcc\_stable} chấp nhận \code{kg\_seed\_concepts} mà
\code{kg\_expand\_node} thực sự sinh ra. Lý do từ chối đổi hẳn --- từ 64/64
\code{no\_compatible\_candidate} sang 28/32 ca đánh giá ứng viên thật, risk 0,046--0,152 ---
nhưng L1 vẫn 0\%, vì tái dùng thẳng đòi \code{query\_norm} giống hệt nên câu diễn đạt lại
luôn rơi vào nhánh vá, mà không slot nào được khai báo.

\textbf{Giai đoạn 2} khai báo slot ở tầng harness. L1 trúng 75\% --- đúng bằng thiết kế bộ dữ
liệu (8 cụm $\times$ 4 biến thể) --- nhưng EM sụp một nửa. Nội dung tái dùng \emph{đúng};
thứ phá điểm là đuôi \code{(Also related: …)} mà \code{\_patch\_response} nối thêm, chứa
\code{token:that}, \code{token:ab} --- các mảnh từ vựng nội bộ dùng để gom bucket, không bao
giờ được thành văn. Hàng rào an toàn không chặn vì \code{\_factual\_projection} \emph{cố ý}
cắt đúng đuôi đó trước khi băm: hợp lý cho phản hồi gia sư, sai cho QA nơi toàn bộ phản hồi
chính là đáp án.

\textbf{Giai đoạn 3} lọc, chỉ giữ \code{concept:}/\code{intent:}. Đây không phải vấn đề riêng
của benchmark: người học cũng sẽ đọc thấy ``(Also related: token:that)''.

\begin{luuy}[Cái L1 đem lại, và hai điều kèm theo]
Ở giai đoạn 3, TRACE-CAG phục vụ \textbf{75\% yêu cầu từ đệm mà chất lượng không đổi}
(EM 50,0\% ở cả hai; F1 lệch 0,5 điểm, trong nhiễu): độ trễ trung bình 1\,735\,ms $\to$
\textbf{397\,ms} (4,4$\times$), p50 479\,ms $\to$ \textbf{5\,ms}, prompt token giảm
\textbf{73\%}.

Đối sánh từng câu với lần chạy không có L1: \textbf{28/32 câu trả lời giống hệt}, số câu đúng
bằng nhau. L1 phát lại trung thành cái mà đường tính đầy đủ sinh ra --- \emph{kể cả đáp án
sai}. Một đáp án lạnh sai sẽ lan ra toàn cụm thay vì chỉ sai một lần.

Chạy lại HotpotQA n=64 với cả ba bản vá cho kết quả \textbf{giống hệt đến từng chỉ số}
(EM 62,5\%, F1 74,7\%, R@5 78,9\%, đệm 48,4\%) và \textbf{L1 vẫn đúng 0\%}. Đó không phải
bản vá thất bại mà là chứng chỉ làm đúng việc: 64 câu hỏi không liên quan, mỗi câu mang ảnh
chụp ngữ cảnh riêng, nên ứng viên chéo câu hỏi bị loại ở \code{mismatch:evidence\_hash}
trước khi chấm risk. L1 kích hoạt ở nơi có câu gần nghĩa và im lặng ở nơi không có.
\end{luuy}

Hai giới hạn phải nêu kèm. Thứ nhất, \code{patchable\_slots} là \textbf{tham số giao thức},
phải công bố cùng chỉ số đệm đúng như ngân sách bằng chứng. Thứ hai, việc khai báo đó nằm
trong \code{model-development/} --- thư mục mà repo LexiLingo bỏ qua, nhưng bản thân nó lại
là một kho git riêng. Thay đổi vẫn theo dõi được ở kho đó, song tính đến 2026-08-30 nó chưa
được commit, nên khi công bố phải nêu mốc commit của \emph{cả hai} kho.

Bản vá cổng ghi và bản lọc token nằm trong \code{api/} nên có hiệu lực ở production. Việc
khai báo slot thì \textbf{chưa}, và không nên bật cho người dùng thật trước khi trả lời được
câu hỏi: một câu hỏi diễn đạt khác có được phép nhận lại câu trả lời đã đệm không? Trong
benchmark thì an toàn nhờ \code{evidence\_hash}; chat production không có mốc cố định tương
đương.

\subsection{Cách biệt so với HippoRAG phần lớn là kỷ luật định dạng}

Đối sánh từng câu với HippoRAG, nối theo chuỗi đáp án vàng duy nhất (52/64 câu; 12 câu bị
loại vì gold trùng lặp \code{yes}/\code{no}/\code{2000}). Nối theo thứ tự dòng là sai --- đã
kiểm chứng: chỉ 1/64 câu khớp gold, đúng như lỗi chỉ số tương đối đã ghi nhận.

Trên 52 câu đó: 24 câu cả hai đúng, 20 câu cả hai sai, \textbf{6 câu chỉ TRACE-CAG đúng, 2
câu chỉ HippoRAG đúng}. Nhưng đọc nội dung 8 câu lệch thì \textbf{6 câu là tranh chấp về
ranh giới span}, chỉ 3 câu là sai kiến thức thật. HippoRAG trả \code{1969-1974} khi vàng là
\code{1969 until 1974}; trả \code{IFFHS World's Best Goalkeeper} khi vàng là
\code{world's best goalkeeper}; một câu trả nguyên một đoạn suy luận. Chiều ngược lại,
TRACE-CAG trả \code{Mumbai, Maharashtra} khi vàng là \code{Mumbai}.

Xác nhận trên toàn tập:

\begin{longtable}{L{4.2cm} R{2.6cm} R{3.4cm} R{2.6cm}}
\toprule
& \textbf{số câu sai EM} & \textbf{F1 $\geq$ 0,5 (gần đúng)} & \textbf{F1 = 0}\\
\midrule
\endhead
HippoRAG~2 & 30 & \textbf{14 (47\%)} & 11\\
TRACE-CAG & 24 & 10 (42\%) & 12\\
\bottomrule
\caption{Gần một nửa số câu HippoRAG mất điểm EM vẫn giữ F1 cao --- đúng dấu hiệu của đáp án
đúng nội dung nhưng sai định dạng.}
\end{longtable}

Độ dài đáp án giải thích phần còn lại: trung vị của cả hai đều là 2 từ, nhưng trung bình của
HippoRAG là \textbf{53,5 từ} (dài nhất 754 từ, 5/64 câu vượt 20 từ) so với 2,3 từ của
TRACE-CAG (dài nhất 12 từ, \textbf{0/64} câu vượt 20 từ). HippoRAG lưỡng cực: hoặc trả lời
ngắn gọn, hoặc xả ra một đoạn suy luận mà EM chấm 0 bất kể nội dung.

\begin{luuy}[Phát biểu đúng về cách biệt với HippoRAG]
Câu ``TRACE-CAG hơn HippoRAG 9,4 điểm EM'' đúng về số nhưng gây hiểu nhầm về nguyên nhân.
Phát biểu chịu được phản biện: \emph{TRACE-CAG có kỷ luật định dạng đáp án tốt hơn ở cùng
mức chất lượng suy luận, với chi phí thấp hơn khoảng 30 lần}. Khoảng cách F1 (5,0 điểm) nhỏ
hơn khoảng cách EM (9,4 điểm) đúng vì lý do này --- F1 tha thứ lỗi định dạng, EM thì không.

Điều này không làm mất giá trị của hệ: trả về span sạch là yêu cầu thật của sản phẩm, và
hậu xử lý trích span của TRACE-CAG đáp ứng được. Nhưng đóng góp đo được nằm ở khâu hậu xử lý
đầu ra, không phải ở kiến trúc đồ thị.
\end{luuy}

\section{Tóm tắt: vì sao mỗi thiết kế tồn tại}

\begin{longtable}{L{4.0cm} L{9.6cm}}
\toprule
\textbf{Quyết định} & \textbf{Lý do}\\
\midrule
\endhead
Chạy song song kg\_expand + diagnose & Hai việc không phụ thuộc dữ liệu của nhau; song song cắt gần một nửa độ trễ nhánh chậm.\\
Ngân sách thời gian theo giai đoạn truy xuất & Trễ có giới hạn trên đoán trước được quan trọng hơn kết quả đầy đủ nhưng không đoán trước được.\\
PedWeight ưu tiên đúng bậc CEFR & Một khái niệm đúng bậc hữu ích hơn nhiều khái niệm sai bậc; tránh làm loãng ngữ cảnh bằng nội dung quá dễ hoặc quá khó.\\
Dự phòng quy tắc khi LLM lỗi & Không bao giờ để pipeline hoàn toàn im lặng, kể cả khi mọi nhà cung cấp LLM đều gặp sự cố.\\
Ba cấu hình thích nghi thay vì một ngưỡng cố định & Đánh đổi tốc độ/chất lượng khác nhau theo áp lực hệ thống và đặc điểm câu hỏi, học trực tuyến không cần huấn luyện ngoại tuyến.\\
IRCoT chỉ kích hoạt có chọn lọc & Suy luận thêm một lượt LLM có chi phí; chỉ đáng làm với câu hỏi thực sự cần bắc cầu qua hai thực thể.\\
Tính lại điểm số ở \code{\_format\_response} & Phản hồi từ đệm và phản hồi tính mới không được phép khác nhau về ngữ nghĩa điểm số.\\
\bottomrule
\caption{Tổng hợp quyết định thiết kế của TRACE-CAG}
\end{longtable}
